\documentclass[10pt]{extarticle}
\usepackage[a4paper,margin=1.5cm]{geometry}
\usepackage[T1]{fontenc}
\usepackage{amsmath,amssymb}
\usepackage{graphicx}
\usepackage[export]{adjustbox}
\usepackage{multicol}
\usepackage{enumitem}
\usepackage{booktabs,array,tabularx}
\setlength{\columnsep}{0.7cm}
\setlength{\parindent}{0pt}
\setlength{\parskip}{3pt}
\pagestyle{empty}
\makeatletter
% figure registry (figdefs.tex is generated by the extract step)
\newcommand{\qfigdef}[4]{\@namedef{qf@file@#1}{#2}\@namedef{qf@w@#1}{#3}\@namedef{qf@wide@#1}{#4}}
\newcommand{\qf@check}[1]{\@ifundefined{qf@file@#1}{\PackageError{score2}{Unknown figure id #1}{}}{}}
\newcommand{\qf@img}[4][]{\includegraphics[width=#2\dimexpr\@nameuse{qf@w@#4}pt\relax,max width=#3,max height=0.45\textheight,#1]{\@nameuse{qf@file@#4}}}
% Full-width questions: finalize wraps every question that contains a wide figure in \qfullbegin...\qfullend,
% so the WHOLE question (stem, figures, options) is set across the page in reading order, then the two columns
% resume. The question's head - its number, stem and first figure (\qfullsplit marks the end of the head) -
% is kept on one page, so a number is never left at a page foot with its figure overleaf; the rest may break
% across pages normally, so no big gaps are left. Without a \qfullsplit, the whole question is kept together
% if it is under 45% of a page. Never write these yourself.
\newif\ifqf@full
\newif\ifqf@boxed
\newsavebox{\qf@box}
\newcommand{\qf@open}{\global\qf@fulltrue\par\global\qf@boxedtrue
  \setbox\qf@box\vbox\bgroup\hsize\textwidth\linewidth\textwidth}
\newcommand{\qf@emit}[1]{\par\egroup\global\qf@boxedfalse
  \ifdim\dimexpr\ht\qf@box+\dp\qf@box\relax<#1\textheight\noindent\usebox{\qf@box}\par\else\unvbox\qf@box\fi}
\newcommand{\qfullsplit}{\ifqf@boxed\qf@emit{0.85}\fi}
\newcommand{\qf@close}{\ifqf@boxed\qf@emit{0.45}\fi\global\qf@fullfalse}
\newcommand{\qfullbegin}{\end{multicols}\qf@open}
\newcommand{\qfullend}{\qf@close\begin{multicols}{2}}
\newcommand{\qfullstart}{\qf@open}
\newcommand{\qfulldone}{\qf@close}
% \qfig[scale]{id}: block figure, centred, directly where it is written (size set by the tool).
% Fallback if a wide figure is not inside a full-width question: break out of the columns in place - no floats.
\newcommand{\qfig}[2][1]{\qf@check{#2}%
  \ifqf@full\qf@blk{#1}{#2}\else
  \ifnum\@nameuse{qf@wide@#2}=1 \expandafter\@firstoftwo\else\expandafter\@secondoftwo\fi
  {\end{multicols}\begin{center}\qf@img{#1}{\textwidth}{#2}\end{center}\begin{multicols}{2}}%
  {\qf@blk{#1}{#2}}\fi}
\newcommand{\qf@blk}[2]{\par\smallskip{\centering\qf@img{#1}{\linewidth}{#2}\par}\smallskip}
% \qfigi[scale]{id}: inline figure (inside option lists / tables / minipages); never breaks columns.
\newcommand{\qfigi}[2][1]{\qf@check{#2}\qf@img[valign=c]{#1}{\linewidth}{#2}}
\makeatother
% structure macros - lint reads these to check completeness and figure placement
\newcommand{\papersec}[1]{\par\bigskip{\centering\Large\bfseries PAPER-#1\par}\medskip}
\newcommand{\partsec}[2]{\par\medskip\noindent{\large\bfseries PART-#1: #2}\par\smallskip}
\newcommand{\sect}[1]{\par\smallskip\noindent\textbf{#1}\par}
\newcommand{\qn}[1]{\par\smallskip\noindent\textbf{#1.}~}
\newcommand{\paraq}[2]{\par\smallskip\noindent\textbf{Paragraph for Questions #1 and #2}\par}
\newenvironment{qoptions}{\begin{itemize}[leftmargin=1.8em,itemsep=1pt,topsep=2pt,label={}]}{\end{itemize}}
\newcommand{\opt}[2]{\item[(#1)] #2}
% figure registry (generated by score2_tools.py finalize)
\qfigdef{p1_phy_s1_q01_1}{Score Advanced - Test 07 (16 Apr 2023)-img/p1_phy_s1_q01_1.png}{168.8}{0}
\qfigdef{p1_phy_s1_q02_1}{Score Advanced - Test 07 (16 Apr 2023)-img/p1_phy_s1_q02_1.png}{179.6}{0}
\qfigdef{p1_phy_s1_q03_1}{Score Advanced - Test 07 (16 Apr 2023)-img/p1_phy_s1_q03_1.png}{148.6}{0}
\qfigdef{p1_phy_s1_q04_1}{Score Advanced - Test 07 (16 Apr 2023)-img/p1_phy_s1_q04_1.png}{162.7}{0}
\qfigdef{p1_phy_s1_q06_1}{Score Advanced - Test 07 (16 Apr 2023)-img/p1_phy_s1_q06_1.png}{104.4}{0}
\qfigdef{p1_phy_s2_para9_10_1}{Score Advanced - Test 07 (16 Apr 2023)-img/p1_phy_s2_para9_10_1.png}{97.7}{0}
\qfigdef{p1_phy_s3_q01_1}{Score Advanced - Test 07 (16 Apr 2023)-img/p1_phy_s3_q01_1.png}{100.4}{0}
\qfigdef{p1_phy_s3_q03_1}{Score Advanced - Test 07 (16 Apr 2023)-img/p1_phy_s3_q03_1.png}{195.9}{0}
\qfigdef{p1_phy_s3_q04_1}{Score Advanced - Test 07 (16 Apr 2023)-img/p1_phy_s3_q04_1.png}{336.5}{1}
\qfigdef{p1_phy_s3_q05_1}{Score Advanced - Test 07 (16 Apr 2023)-img/p1_phy_s3_q05_1.png}{78.2}{0}
\qfigdef{p1_phy_s3_q06_1}{Score Advanced - Test 07 (16 Apr 2023)-img/p1_phy_s3_q06_1.png}{154.6}{0}
\qfigdef{p1_chem_s1_q02_1}{Score Advanced - Test 07 (16 Apr 2023)-img/p1_chem_s1_q02_1.png}{126.4}{0}
\qfigdef{p1_chem_s1_q04_A}{Score Advanced - Test 07 (16 Apr 2023)-img/p1_chem_s1_q04_A.png}{52.5}{0}
\qfigdef{p1_chem_s1_q04_B}{Score Advanced - Test 07 (16 Apr 2023)-img/p1_chem_s1_q04_B.png}{41.5}{0}
\qfigdef{p1_chem_s1_q04_C}{Score Advanced - Test 07 (16 Apr 2023)-img/p1_chem_s1_q04_C.png}{63.5}{0}
\qfigdef{p1_chem_s1_q04_D}{Score Advanced - Test 07 (16 Apr 2023)-img/p1_chem_s1_q04_D.png}{77.9}{0}
\qfigdef{p1_chem_s1_q05_A}{Score Advanced - Test 07 (16 Apr 2023)-img/p1_chem_s1_q05_A.png}{169.4}{0}
\qfigdef{p1_chem_s1_q05_B}{Score Advanced - Test 07 (16 Apr 2023)-img/p1_chem_s1_q05_B.png}{172.0}{0}
\qfigdef{p1_chem_s1_q05_C}{Score Advanced - Test 07 (16 Apr 2023)-img/p1_chem_s1_q05_C.png}{205.8}{0}
\qfigdef{p1_chem_s1_q05_D}{Score Advanced - Test 07 (16 Apr 2023)-img/p1_chem_s1_q05_D.png}{213.1}{0}
\qfigdef{p1_chem_s1_q06_1}{Score Advanced - Test 07 (16 Apr 2023)-img/p1_chem_s1_q06_1.png}{145.7}{0}
\qfigdef{p1_chem_s2_q07_A}{Score Advanced - Test 07 (16 Apr 2023)-img/p1_chem_s2_q07_A.png}{73.5}{0}
\qfigdef{p1_chem_s2_q07_B}{Score Advanced - Test 07 (16 Apr 2023)-img/p1_chem_s2_q07_B.png}{67.7}{0}
\qfigdef{p1_chem_s2_q07_C}{Score Advanced - Test 07 (16 Apr 2023)-img/p1_chem_s2_q07_C.png}{95.6}{0}
\qfigdef{p1_chem_s2_q07_D}{Score Advanced - Test 07 (16 Apr 2023)-img/p1_chem_s2_q07_D.png}{93.1}{0}
\qfigdef{p1_chem_s2_q08_A}{Score Advanced - Test 07 (16 Apr 2023)-img/p1_chem_s2_q08_A.png}{40.2}{0}
\qfigdef{p1_chem_s2_q08_B}{Score Advanced - Test 07 (16 Apr 2023)-img/p1_chem_s2_q08_B.png}{54.4}{0}
\qfigdef{p1_chem_s2_q08_C}{Score Advanced - Test 07 (16 Apr 2023)-img/p1_chem_s2_q08_C.png}{43.1}{0}
\qfigdef{p1_chem_s2_q08_D}{Score Advanced - Test 07 (16 Apr 2023)-img/p1_chem_s2_q08_D.png}{43.0}{0}
\qfigdef{p1_chem_s2_para9_10_1}{Score Advanced - Test 07 (16 Apr 2023)-img/p1_chem_s2_para9_10_1.png}{108.9}{0}
\qfigdef{p1_chem_s3_q02_1}{Score Advanced - Test 07 (16 Apr 2023)-img/p1_chem_s3_q02_1.png}{46.8}{0}
\qfigdef{p1_chem_s3_q02_2}{Score Advanced - Test 07 (16 Apr 2023)-img/p1_chem_s3_q02_2.png}{311.0}{1}
\qfigdef{p1_chem_s3_q02_3}{Score Advanced - Test 07 (16 Apr 2023)-img/p1_chem_s3_q02_3.png}{291.2}{1}
\qfigdef{p1_chem_s3_q02_4}{Score Advanced - Test 07 (16 Apr 2023)-img/p1_chem_s3_q02_4.png}{144.2}{0}
\qfigdef{p1_chem_s3_q05_1}{Score Advanced - Test 07 (16 Apr 2023)-img/p1_chem_s3_q05_1.png}{92.0}{0}
\qfigdef{p2_phy_s1_q01_1}{Score Advanced - Test 07 (16 Apr 2023)-img/p2_phy_s1_q01_1.png}{97.3}{0}
\qfigdef{p2_phy_s1_q02_1}{Score Advanced - Test 07 (16 Apr 2023)-img/p2_phy_s1_q02_1.png}{108.2}{0}
\qfigdef{p2_phy_s1_q04_1}{Score Advanced - Test 07 (16 Apr 2023)-img/p2_phy_s1_q04_1.png}{73.0}{0}
\qfigdef{p2_phy_s1_q05_1}{Score Advanced - Test 07 (16 Apr 2023)-img/p2_phy_s1_q05_1.png}{63.5}{0}
\qfigdef{p2_phy_s2_q07_1}{Score Advanced - Test 07 (16 Apr 2023)-img/p2_phy_s2_q07_1.png}{97.7}{0}
\qfigdef{p2_phy_s2_q07_2}{Score Advanced - Test 07 (16 Apr 2023)-img/p2_phy_s2_q07_2.png}{95.4}{0}
\qfigdef{p2_phy_s2_q07_3}{Score Advanced - Test 07 (16 Apr 2023)-img/p2_phy_s2_q07_3.png}{95.4}{0}
\qfigdef{p2_phy_s2_q07_4}{Score Advanced - Test 07 (16 Apr 2023)-img/p2_phy_s2_q07_4.png}{124.4}{0}
\qfigdef{p2_phy_s2_q08_1}{Score Advanced - Test 07 (16 Apr 2023)-img/p2_phy_s2_q08_1.png}{79.4}{0}
\qfigdef{p2_phy_s2_q09_1}{Score Advanced - Test 07 (16 Apr 2023)-img/p2_phy_s2_q09_1.png}{161.3}{0}
\qfigdef{p2_phy_s2_q10_1}{Score Advanced - Test 07 (16 Apr 2023)-img/p2_phy_s2_q10_1.png}{187.5}{0}
\qfigdef{p2_phy_s3_q01_1}{Score Advanced - Test 07 (16 Apr 2023)-img/p2_phy_s3_q01_1.png}{105.3}{0}
\qfigdef{p2_phy_s3_q02_1}{Score Advanced - Test 07 (16 Apr 2023)-img/p2_phy_s3_q02_1.png}{169.7}{0}
\qfigdef{p2_phy_s3_q04_1}{Score Advanced - Test 07 (16 Apr 2023)-img/p2_phy_s3_q04_1.png}{182.9}{0}
\qfigdef{p2_phy_s3_q04_2}{Score Advanced - Test 07 (16 Apr 2023)-img/p2_phy_s3_q04_2.png}{82.3}{0}
\qfigdef{p2_phy_s3_q06_1}{Score Advanced - Test 07 (16 Apr 2023)-img/p2_phy_s3_q06_1.png}{176.5}{0}
\qfigdef{p2_phy_s3_q07_1}{Score Advanced - Test 07 (16 Apr 2023)-img/p2_phy_s3_q07_1.png}{190.6}{0}
\qfigdef{p2_phy_s3_q08_1}{Score Advanced - Test 07 (16 Apr 2023)-img/p2_phy_s3_q08_1.png}{68.0}{0}
\qfigdef{p2_chem_s1_q03_A}{Score Advanced - Test 07 (16 Apr 2023)-img/p2_chem_s1_q03_A.png}{73.5}{0}
\qfigdef{p2_chem_s1_q03_B}{Score Advanced - Test 07 (16 Apr 2023)-img/p2_chem_s1_q03_B.png}{115.9}{0}
\qfigdef{p2_chem_s1_q03_C}{Score Advanced - Test 07 (16 Apr 2023)-img/p2_chem_s1_q03_C.png}{109.0}{0}
\qfigdef{p2_chem_s1_q03_D}{Score Advanced - Test 07 (16 Apr 2023)-img/p2_chem_s1_q03_D.png}{103.5}{0}
\qfigdef{p2_chem_s1_q06_1}{Score Advanced - Test 07 (16 Apr 2023)-img/p2_chem_s1_q06_1.png}{107.5}{0}
\qfigdef{p2_chem_s1_q06_2}{Score Advanced - Test 07 (16 Apr 2023)-img/p2_chem_s1_q06_2.png}{48.2}{0}
\qfigdef{p2_chem_s1_q06_3}{Score Advanced - Test 07 (16 Apr 2023)-img/p2_chem_s1_q06_3.png}{44.9}{0}
\qfigdef{p2_chem_s1_q06_4}{Score Advanced - Test 07 (16 Apr 2023)-img/p2_chem_s1_q06_4.png}{170.0}{0}
\qfigdef{p2_chem_s1_q06_5}{Score Advanced - Test 07 (16 Apr 2023)-img/p2_chem_s1_q06_5.png}{47.0}{0}
\qfigdef{p2_chem_s1_q06_6}{Score Advanced - Test 07 (16 Apr 2023)-img/p2_chem_s1_q06_6.png}{91.8}{0}
\qfigdef{p2_chem_s1_q06_7}{Score Advanced - Test 07 (16 Apr 2023)-img/p2_chem_s1_q06_7.png}{215.0}{0}
\qfigdef{p2_chem_s2_q07_4}{Score Advanced - Test 07 (16 Apr 2023)-img/p2_chem_s2_q07_4.png}{148.7}{0}
\qfigdef{p2_chem_s2_q07_5}{Score Advanced - Test 07 (16 Apr 2023)-img/p2_chem_s2_q07_5.png}{138.1}{0}
\qfigdef{p2_chem_s2_q09_1}{Score Advanced - Test 07 (16 Apr 2023)-img/p2_chem_s2_q09_1.png}{66.4}{0}
\qfigdef{p2_chem_s2_q09_2}{Score Advanced - Test 07 (16 Apr 2023)-img/p2_chem_s2_q09_2.png}{79.4}{0}
\qfigdef{p2_chem_s2_q09_3}{Score Advanced - Test 07 (16 Apr 2023)-img/p2_chem_s2_q09_3.png}{68.6}{0}
\qfigdef{p2_chem_s2_q09_4}{Score Advanced - Test 07 (16 Apr 2023)-img/p2_chem_s2_q09_4.png}{63.0}{0}
\qfigdef{p2_chem_s3_q03_1}{Score Advanced - Test 07 (16 Apr 2023)-img/p2_chem_s3_q03_1.png}{205.4}{0}
\qfigdef{p2_chem_s3_q06_A}{Score Advanced - Test 07 (16 Apr 2023)-img/p2_chem_s3_q06_A.png}{54.1}{0}
\qfigdef{p2_chem_s3_q06_B}{Score Advanced - Test 07 (16 Apr 2023)-img/p2_chem_s3_q06_B.png}{54.1}{0}
\qfigdef{p2_chem_s3_q06_C}{Score Advanced - Test 07 (16 Apr 2023)-img/p2_chem_s3_q06_C.png}{61.5}{0}
\qfigdef{p2_chem_s3_q06_D}{Score Advanced - Test 07 (16 Apr 2023)-img/p2_chem_s3_q06_D.png}{115.3}{0}
\qfigdef{p2_chem_s3_q06_5}{Score Advanced - Test 07 (16 Apr 2023)-img/p2_chem_s3_q06_5.png}{201.3}{0}
\qfigdef{p2_chem_s3_q06_6}{Score Advanced - Test 07 (16 Apr 2023)-img/p2_chem_s3_q06_6.png}{200.8}{0}
\raggedcolumns
\makeatletter\@beginparpenalty=10000\makeatother
\newcolumntype{L}[1]{>{\raggedright\arraybackslash}p{#1}}
\begin{document}
\begin{center}{\Large\bfseries Score Advanced --- Test 07 (16 Apr 2023)}\end{center}
\begin{multicols}{2}
\raggedcolumns
\papersec{1}
\partsec{1}{PHYSICS}
\sect{SECTION-I (i)}
\qn{1} Two infinite current carrying wires lie overlapping each other along x \& y axes carrying current $\text{I}_2 = \text{I}_0\cos\!\left(\omega\text{t}+\dfrac{\pi}{4}\right)$ and $\text{I}_1 = 2\text{I}_0\sin\!\left(\omega\text{t}+\dfrac{\pi}{4}\right)$. A rectangular frame of negligible resistance with capacitor of capacitance C is placed in the xy plane as shown. Select correct options.\par
\qfig{p1_phy_s1_q01_1}
\begin{qoptions}
\opt{A}{Maximum emf induced in frame is $\dfrac{\mu_0\text{I}_0\omega\ell\text{n}2}{2\pi}$}
\opt{B}{Maximum emf induced in frame is $\dfrac{\sqrt{2}\mu_0\text{I}_0\omega\ell\text{n}2}{\pi}$}
\opt{C}{Maximum charge on capacitor is $\dfrac{\sqrt{2}\mu_0\text{I}_0\omega\text{C}\ell\text{n}2}{\pi}$}
\opt{D}{Current in frame at t = $\dfrac{\pi}{\omega}$ is zero}
\end{qoptions}
\qn{2} In the YDSE setup shown in figure a monochromatic of light of wavelength $\lambda$ is incident on slit S. Distance D is very large as compared to separation d between slits S$_1$ and S$_2$. Choose the correct alternative(s):\par
\qfig{p1_phy_s1_q02_1}
\begin{qoptions}
\opt{A}{Minimum value of d for which there is a bright fringe at O is $\sqrt{\dfrac{2\lambda\text{D}}{3}}$}
\opt{B}{Minimum value of d for which there is a bright fringe at O is $\sqrt{\dfrac{3\lambda\text{D}}{2}}$}
\opt{C}{Position of central bright fringe for minimum value of d is $\sqrt{\dfrac{2\lambda\text{D}}{3}}$ above O}
\opt{D}{Position of central bright fringe for minimum value of d is $\sqrt{6\lambda\text{D}}$ above O}
\end{qoptions}
\qn{3} A charged particle (+q) is moving simple harmonically on the x-axis with its mean position at origin. Amplitude of the particle is A and its angular frequency is $\omega$. Then choose the correct option(s) :\par
\qfig{p1_phy_s1_q03_1}
\begin{qoptions}
\opt{A}{The magnitude of magnetic field at (2A, 0) will change periodically with period 2$\pi/\omega$.}
\opt{B}{The maximum magnitude of the magnetic field at (0, A) is $\dfrac{\mu_0}{4\pi}\cdot\dfrac{q\omega}{A}$}
\opt{C}{The magnetic field at (A, A) at the moment the particle passes through $\left(\dfrac{A}{2},0\right)$, will be $\dfrac{\sqrt{3}\cdot\mu_0 q\omega}{5\sqrt{5}\pi A}$}
\opt{D}{The magnitudes of magnetic field at (0, A) and (0, --A) will be same at any time.}
\end{qoptions}
\qn{4} An ideal liquid of density $\rho$ is filled in a cylindrical container upto (3/4)$^{\text{th}}$ height. Now the liquid is being rotated about vertical axis passing through its axis of symmetry with constant angular velocity $\omega$, such that liquid is on the verge of falling out of the container and free surface of liquid forms a paraboloid. Whole situation is shown in the figure. Choose the correct option(s).\par
\qfig{p1_phy_s1_q04_1}
\begin{qoptions}
\opt{A}{h$_0$ = $\dfrac{\text{h}}{4}$}
\opt{B}{$\omega$ = $\sqrt{\dfrac{\text{gh}}{\text{R}^2}}$}
\opt{C}{$\left|\dfrac{\text{dp}}{\text{dx}}\right| = \rho\omega^2\text{x}$}
\opt{D}{$\left|\dfrac{\text{dp}}{\text{dy}}\right| = \rho\text{g}$}
\end{qoptions}
\qn{5} The radii of a spherical capacitor are equal to a and b (b $>$ a). The space between them is filled with a dielectric of dielectric constant K and resistivity $\rho$. At t = 0, the inner electrode is given a charge q$_0$. Choose the correct option(s) :
\begin{qoptions}
\opt{A}{Charge q on the inner electrode as a function of time is given by q = q$_0$ $\exp\!\left(-\dfrac{\text{t}}{\rho\text{K}\varepsilon_0}\right)$}
\opt{B}{In a short time, the charge on the inner electrode will become zero}
\opt{C}{After a long time, the charge on the outer sphere will become q$_0$}
\opt{D}{The total amount of heat generated during the spreading of charge will be given by H = $\left(\dfrac{1}{\text{a}}-\dfrac{1}{\text{b}}\right)\dfrac{\text{q}_0^2}{8\pi\varepsilon_0\text{K}}$}
\end{qoptions}
\qn{6} A small object moves counter clockwise along the circular path whose centre is at origin as shown in figure. As it moves along the path, its acceleration vector continuously points towards point S. Then the object\par
\qfig{p1_phy_s1_q06_1}
\begin{qoptions}
\opt{A}{Speed up as it moves from A to C via B.}
\opt{B}{Slows down as it moves from A to C via B.}
\opt{C}{Slows down as it moves from C to A via D.}
\opt{D}{Speed up as it moves from C to A via D.}
\end{qoptions}
\sect{SECTION-I (ii)}
\paraq{7}{8}
If a plane simple harmonic wave is traveling towards positive x axis, its wave function is given by\par
S =S$_0$ sin ($\omega$t -- kx)\par
Here S is longitudinal displacement of any particle. If we take a plane parallel to yz plane all particles in this plane have same phase. This plane is called wavefront.\par
If a plane wave is traveling along direction $\hat{n}$, its wave function can be written as :\par
S = S$_0$ sin ($\omega$t -- k$\hat{n}\cdot\vec{r}$)\par
Let's call k$\hat{n}$ = $\vec{k}$ (a wave vector)\par
S = S$_0$ sin($\omega$t -- $\vec{k}\cdot\vec{r}$)\par
Wave vector points in direction of phase velocity. In other words, the wave vector points in the normal direction to the wavefront.\par
\qn{7} Consider a wave function :\par
S = S$_0$ sin($\omega$t -- 2x + 2y + z) Here all variables are in SI units.\par
Select correct Statement :
\begin{qoptions}
\opt{A}{Wavelength of the wave is $\dfrac{2\pi}{3}$.}
\opt{B}{Particle (0, 0, 0) and (1, 1, 2) are on same wavefront.}
\opt{C}{The wave vector is at angle $\cos^{-1}\!\left(\dfrac{1}{3}\right)$ from x-axis.}
\opt{D}{The wave vector is at equal angle from positive x-axis and negative z-axis.}
\end{qoptions}
\qn{8} Consider following time wave function.\par
S$_1$ = S$_0$ sin($\omega$t -- 2x + 2y + z)\par
S$_2$ = S$_0$ sin($\omega$t + 2x -- 2y -- z)\par
An observer is traveling towards +x axis with speed $\pi$ m/s. The beat frequency observed by the observer is :
\begin{qoptions}
\opt{A}{$\dfrac{1}{2}$}
\opt{B}{2}
\opt{C}{1}
\opt{D}{$\dfrac{3}{2}$}
\end{qoptions}
\paraq{9}{10}
A long rectangular slab of transparent medium of thickness d is placed on table with length parallel to x-axis and width parallel to the y-axis. A ray of light traveling nearly along y-axis at origin. The refractive index $\mu$ of the medium varies as $\mu = \sqrt{1+\text{e}^{\text{x/d}}}$. The refractive index of the air is 1.\par
\qfig{p1_phy_s2_para9_10_1}
\qn{9} Determine the x-coordinate of the point A, where the ray intersect the upper surface of the slab-air boundary.
\begin{qoptions}
\opt{A}{d(ln2)}
\opt{B}{2d(ln2)}
\opt{C}{$\dfrac{\text{d(ln 2)}}{2}$}
\opt{D}{None of these}
\end{qoptions}
\qn{10} Write down the refractive index of the medium at A.
\begin{qoptions}
\opt{A}{$\sqrt{5}$}
\opt{B}{$\sqrt{6}$}
\opt{C}{$\sqrt{7}$}
\opt{D}{$\sqrt{8}$}
\end{qoptions}
\sect{SECTION-II}
\qn{1} Four short electric dipoles, 3 each of dipole moment P and the fourth of dipole moment 2P are fixed at the vertices of an square of side length $\sqrt{2}$r. The dipole moments are oriented as shown in diagram and the centre of the square lies at the origin. The electric field strength vector at the origin is given as $\dfrac{\text{KP}}{\text{r}^3}(\text{x}\hat{\text{i}}+\text{y}\hat{\text{j}})$, then find value of (x + y).\par
\qfig{p1_phy_s3_q01_1}
\qn{2} An inductor (L), a resistance (R) and a ideal battery (emf $\varepsilon$) are connected in series. Circuit is switched on at t = 0. The ratio of energy stored in the inductor at time t = L/R and energy dissipated in resistance up to that instant is $\dfrac{\text{U}_1}{\text{U}_2} = \dfrac{1}{\alpha}\,\dfrac{\left(1-\frac{1}{\text{e}}\right)^2}{\left(\frac{\beta}{\text{e}}-\frac{1}{\gamma\text{e}^2}-\frac{1}{\delta}\right)}$. Here $\alpha$, $\beta$, $\gamma$ and $\delta$ are integers, e is natural logarithm. Find ($\alpha$ + $\beta$ + $\gamma$ + $\delta$).
\qn{3} Consider the shown diagram where the two chambers separated by piston-spring arrangement contain equal amounts of same ideal gas. Initially when the temperatures of the gas in both the chambers are kept at 300 K, the compression in the spring is 10 cm. The temperature of the left and the right chambers are now raised slowly to 400 K and 500 K respectively. The pistons are free to slide (i.e. no friction between piston and chamber). Final compression in the spring is found to be $\dfrac{\alpha}{60}$m. Find $\alpha$. (Use $\sqrt{481}$ = 22 )\par
\qfig{p1_phy_s3_q03_1}
\qfullbegin
\qn{4} A uniform heavy plank of length L is resting on two rods A and B of identical length and same area of cross section but made up of different materials. It is observed that to keep the plank horizontal, one rod is to be at the end of the plank (say rod A) and another rod (say rod B) is to be at separation $\dfrac{\text{L}}{4}$ from centre of plank as shown in figure. Now the rods are kept between two rigid walls such that separation between the walls is equal to the natural length of the rod. Now temperature of both the rods is gradually increased in same manner, it is observed that thermal stress developed in rod A is always thrice the thermal stress developed in rod B. Ratio of coefficient of linear expansion of rod A to coefficient of linear expansion of rod B will be :\par
\qfig{p1_phy_s3_q04_1}
\qfullend
\qn{5} The chain of length L is released from rest on a smooth fixed incline with x = 0 as shown in the figure. Determine velocity v of the chain when a half of the length has fallen of incline (Take L = $\dfrac{16}{25}$m) in m/s. (Neglect edge effect of inclined). (Take g = 10 m/s$^2$)\par
\qfig{p1_phy_s3_q05_1}
\qn{6} A right angled prism is placed with air all around it. A horizontal incident ray strikes the face PQ of prism normally, then undergoes total internal reflection twice as shown and exists the prism from face PR horizontally. The refractive index of prism is $\mu$ and angle $\theta$ = 15º. If the value of $\mu$ is $\sqrt{\text{n}}+\dfrac{1}{2}$ then find n.\par
\qfig{p1_phy_s3_q06_1}
\qn{7} Starting from t = 0, activity of a radioactive sample is A$_1$ at time t$_1$ and A$_2$ at time t$_2$. If $\tau$ is average life of sample then the number of nuclei decayed in time (t$_2$ -- t$_1$) is $\beta$(A$_1$ -- A$_2$) in proper unit. Here $\beta$ is an integer. Find $\beta$. (Given that t$_1$ = 5 sec, t$_2$ = 8 sec and $\tau$ = 7 sec)
\qn{8} In a screw gauge whose least count is 0.005 mm, when no object is present and the screw gauge gap is completely closed the 4$^{\text{th}}$ division on its circular scale coincides with the reference line on main scale reading zero. Now a wire is placed in it to measure its diameter, the reading on the main scale advances by 6 divisions, whereas the circular scale reading becomes 3 times to the corresponding reading when no object was placed. If there are 200 divisions on circular scale, then find the diameter of wire in mm.
\partsec{2}{CHEMISTRY}
\sect{SECTION-I (i)}
\qn{1} A, B and C are three complexes of chromium (III) with empirical formula H$_{12}$O$_6$Cl$_3$Cr. All the three complexes have water and chloride ion as ligands. Complex A does not react with concentrated H$_2$SO$_4$, whereas complexes B and C lose 6.75\% and 13.5\% of their original mass, respectively, on treatment with concentrated H$_2$SO$_4$. The compound A,B and C respectively are
\begin{qoptions}
\opt{A}{A = [Cr(HO)$_5$Cl] Cl$_2\cdot$H$_2$O,B = [Cr(H$_2$O)$_4$] Cl$_3$, C = [Cr(H$_2$O)$_6$Cl$_2$]}
\opt{B}{A = [Cr(H$_2$O)$_6$] Cl$_3$, B = [Cr(H$_2$O)$_5$Cl] Cl$_2\cdot$H$_2$O, C = [Cr(H$_2$O)$_4$Cl$_2$]Cl$\cdot$2H$_2$O}
\opt{C}{A = [Cr(H$_2$O)$_5$] Cl$_3\cdot$H$_2$O, B = [Cr(H$_2$O)$_6$] Cl$_3$, C = [Cr(H$_2$O)$_4$ Cl$_2$] Cl$\cdot$H$_2$O}
\opt{D}{A = [Cr(H$_2$O)$_6$ ] Cl$_3$, B = [Cr(H$_2$O)$_4$ Cl$_2$] Cl$\cdot$2H$_2$O, C = [Cr(H$_2$O)$_5$Cl] Cl$_2\cdot$H$_2$O}
\end{qoptions}
\qn{2} Attainment of the equilibrium A(g) $\rightleftharpoons$ 2C(g) + B(g) gave the following of graph. Find the correct option(s).\par
\qfig{p1_chem_s1_q02_1}
\begin{qoptions}
\opt{A}{Equilibrium const. K$_\text{C}$ for above reaction is 128 M$^2$.}
\opt{B}{At equilibrium \% dissociation of A is 40\%}
\opt{C}{At equilibrium \% dissociation of A is 30\%}
\opt{D}{None of these}
\end{qoptions}
\qn{3} Which of the following statement(s) is/are correct?
\begin{qoptions}
\opt{A}{The degree of hydration decreases on moving down the group from Li$^+$ to Cs$^+$}
\opt{B}{Li$_2$O$_2$ is prepared by adding hydrogen peroxide and alcohol to a concentrated solution of lithium hydroxide}
\opt{C}{Lattice energy of ionic solids having the same anion decreases with increase in the size of the cation}
\opt{D}{Lithium nitrate evolves nitrogen monoxide and oxygen on heating}
\end{qoptions}
\qn{4} A compound `X' evolves CO$_2$ when treated with NaHCO$_3$ but gives negative test with tollen's reagent and fehling's solution. On treatment with CHCI$_3$, KOH, afoul smelling organic compound is formed. If `X' do not give positive iodoform test, which of the following satisfy the conditions to be `X'.
\begin{qoptions}
\opt{A}{\qfigi{p1_chem_s1_q04_A}}
\opt{B}{\qfigi{p1_chem_s1_q04_B}}
\opt{C}{\qfigi{p1_chem_s1_q04_C}}
\opt{D}{\qfigi{p1_chem_s1_q04_D}}
\end{qoptions}
\qn{5} Among the following, which reaction will give N-ethylcyclopentylamine as major product?
\begin{qoptions}
\opt{A}{\qfigi{p1_chem_s1_q05_A}}
\opt{B}{\qfigi{p1_chem_s1_q05_B}}
\opt{C}{\qfigi{p1_chem_s1_q05_C}}
\opt{D}{\qfigi{p1_chem_s1_q05_D}}
\end{qoptions}
\qn{6} Consider the following reversible processes and choose the correct answer(s) from the options given below:\par
\qfig{p1_chem_s1_q06_1}
\begin{qoptions}
\opt{A}{$\Delta$S$_{\text{x}\to\text{y}}$ = 95.74 JK$^{-1}$}
\opt{B}{$\Delta$S$_{\text{x}\to\text{z}}$ = 0}
\opt{C}{$\Delta$S$_{\text{x}\to\text{y}}$ = $\Delta$S$_{\text{x}\to\text{z}}$ + $\Delta$S$_{\text{z}\to\text{y}}$ = 95.74 JK$^{-1}$}
\opt{D}{$\Delta$S$_{\text{x}\to\text{z}\to\text{y}}$ = 0}
\end{qoptions}
\sect{SECTION-I (ii)}
\paraq{7}{8}
An organic compound A, C$_8$H$_4$O$_3$ , in dry benzene in the presence of anhydrous AlCl$_3$, gives compound B. The compound B on treatment with PCl$_5$ followed by reaction withH$_2$/Pd, BaSO$_4$ gives compound C, which on reaction with hydrazine gives a cyclised compound D(C$_{14}$H$_{10}$N$_2$).\par
\qn{7} The compound `B' is
\begin{qoptions}
\opt{A}{\qfigi{p1_chem_s2_q07_A}}
\opt{B}{\qfigi{p1_chem_s2_q07_B}}
\opt{C}{\qfigi{p1_chem_s2_q07_C}}
\opt{D}{\qfigi{p1_chem_s2_q07_D}}
\end{qoptions}
\qn{8} The compound `D' is
\begin{qoptions}
\opt{A}{\qfigi{p1_chem_s2_q08_A}}
\opt{B}{\qfigi{p1_chem_s2_q08_B}}
\opt{C}{\qfigi{p1_chem_s2_q08_C}}
\opt{D}{\qfigi{p1_chem_s2_q08_D}}
\end{qoptions}
\paraq{9}{10}
Two moles of an ideal monoatomic gas is taken along the path ABC as shown in following figure. Given T$_\text{B}$ = 2T$_\text{A}$\par
\qfig{p1_chem_s2_para9_10_1}
\qn{9} Increase in internal energy in AB process is
\begin{qoptions}
\opt{A}{3RT$_\text{A}$}
\opt{B}{$\dfrac{9}{4}$RT$_\text{A}$}
\opt{C}{$\dfrac{3}{4}$RT$_\text{B}$}
\opt{D}{$\dfrac{3}{2}$RT$_\text{A}$}
\end{qoptions}
\qn{10} Work done by the gas in process ABC
\begin{qoptions}
\opt{A}{2RT$_\text{A}$}
\opt{B}{--2RT$_\text{A}$}
\opt{C}{4RT$_\text{A}$}
\opt{D}{None of these}
\end{qoptions}
\sect{SECTION-II}
\qn{1} Among the given molecules, the number of metal sulphides precipitated by passing H$_2$S in acidic medium are ..........\par
ZnS, PbS, CuS, NiS, MnS, CdS, CoS, HgS, SnS
\qfullbegin
\qn{2} X = number of reaction in which \qfigi{p1_chem_s3_q02_1} will form as final product.\par
\qfig{p1_chem_s3_q02_2}
\qfullsplit
\qfig{p1_chem_s3_q02_3}
\qfig{p1_chem_s3_q02_4}
Find value of `X' ?
\qfullend
\qn{3} A cube shaped crystal of an alkali metal, 2 mm on an edge, was vaporised in a 500.00 mL flask. The pressure of the resulting vapour was 6.08 mm at 727$^\circ$ C. The structure of the solid metal is known to be body centred cubic. If volume of each unit cell of this crystal is [P $\times$ 10$^{-28}$ m$^3$] then find value of `P'?\par
[Given : R = 0.08 L-atm mol$^{-1}$K$^{-1}$ , N$_\text{A}$ = 6 $\times$ 10$^{23}$]
\qn{4} The ionisation constant of benzoic acid is 6.46 $\times$ 10$^{-6}$ and K$_\text{sp}$ for silver benzoate is 2.5 $\times$ 10$^{-13}$. The number of times silver benzoate is more soluble in a buffer of pH 3.19 compared to its solubility in pure water will be .........\par
[Given : Antilog [--3.19] = 6.46 $\times$ 10$^{-4}$ M]
\qn{5} X = total number of $\beta$-elimination products. Find the value of X ?\par
\qfig{p1_chem_s3_q05_1}
\qn{6} A mixture of ethyl alcohol and propyl alcohol has a vapour pressure of 320 torr at 300 K. Calculate moles of ethyl alcohol in vapour phase if total moles of mixture in vapour phase are 10.29.\par
Given : Vapour pressure of pure ethyl alcohol = 350 torr ; Vapour pressure of pure propyl alcohol = 200 torr.
\qn{7} 40 ml of 0.5 M aqueous solution of oxalic acid (molar mass = 126) is shaken with 25.2 g of wood charcoal. The final concentration of the solution after adsorption of oxalic acid is 0.25 M. If the weight of oxalic acid adsorbed per g of carbon is represented as y $\times$ 10$^{-2}$ then find the value of y. (Assume volume of oxalic acid remains constant)
\qn{8} Two solid compounds A and B dissociates into gaseous products at 293 K as\par
(i) A(s) $\rightleftharpoons$ X(g) + H$_2$S(g).\par
(ii) B (s) $\rightleftharpoons$ Y(g) + H$_2$S(g).\par
Under the given condition, pressure over excess solid A when taken alone in a vessel is 68 mm and that over excess solid B when taken alone in a vessel is 50 mm. The total pressure [in mm Hg] if decomposition of both solids A \& B takes place in the same vessel is ......... .
\partsec{3}{MATHEMATICS}
\sect{SECTION-I (i)}
\qn{1} Let J = $\displaystyle\int_{-2}^{2}\left(\cot^{-1}\dfrac{1}{\text{x}}+\cot^{-1}\text{x}\right)\text{dx}$ and K = $\displaystyle\int_{-2\pi}^{7\pi}\dfrac{\sin\text{x}}{|\sin\text{x}|}\text{dx}$. Then which of the following alternative(s) is/are correct?
\begin{qoptions}
\opt{A}{2J + 3K = 11$\pi$}
\opt{B}{4J$^2$ + K$^2$ = 26$\pi^2$}
\opt{C}{J is period of $\left(\tan\dfrac{\text{x}}{8}\right)^4+\left(\cot\dfrac{\text{x}}{8}\right)^4$}
\opt{D}{$\dfrac{\text{J}}{\text{K}}$ = 4}
\end{qoptions}
\qn{2} A fair dice is rolled four times. Find the probability that each number is not greater than the preceding number. When it is known that it shows `1' only once \& that too in the fourth round.
\begin{qoptions}
\opt{A}{$\dfrac{7}{25}$}
\opt{B}{less than $\dfrac{1}{3}$}
\opt{C}{greater than $\dfrac{1}{3}$}
\opt{D}{$\dfrac{35}{216}$}
\end{qoptions}
\qn{3} Let AB be a chord of the parabola x$^2$ = 4y. A circle drawn with AB as diameter passes through the vertex C of the parabola. If area of $\Delta$ABC = 20 sq.units then the coordinates of A can be
\begin{qoptions}
\opt{A}{(6, 9)}
\opt{B}{(--6, 9)}
\opt{C}{(--8, 16)}
\opt{D}{(8, 16)}
\end{qoptions}
\qn{4} For a 3 $\times$ 3 invertible matrix A satisfying the characteristic equation A$^3$ + pA$^2$ + qA -- rI = 0, which of the following is/are true [Here tr(A) = trace of matrix A, det(A)=determinant value of matrix A]
\begin{qoptions}
\opt{A}{p = --tr(A)}
\opt{B}{q = $\dfrac{(\text{tr}(\text{A}))^2-\text{tr}(\text{A}^2)}{2}$}
\opt{C}{r = det(A)}
\opt{D}{A$^{-1}$ = $\dfrac{-2\,(\text{A}-\text{tr}(\text{A})\text{I}_3)}{(\text{tr}(\text{A}))^2-\text{tr}(\text{A}^2)}$}
\end{qoptions}
\qn{5} Given that a$_1$,a$_2$,a$_3$ is an A.P in that order, a$_1$ + a$_2$ + a$_3$ = 15; b$_1$,b$_2$,b$_3$ is a G.P. in that order, and b$_1$b$_2$b$_3$ = 27. If a$_1$ + b$_1$,a$_2$ + b$_2$,a$_3$ + b$_3$ are positive integers and form a G.P. in that order, then
\begin{qoptions}
\opt{A}{maximum possible value of a$_3$ is 15 (a$_1$,a$_2$,a$_3$ are positive integer and unequal)}
\opt{B}{minimum possible value of a$_3$ is 3 (a$_1$,a$_2$,a$_3$ are positive integer and unequal)}
\opt{C}{maximum possible value of a$_3$ is 9 + $\dfrac{55+7\sqrt{61}}{2}$}
\opt{D}{maximum possible value of a$_3$ is 4 + $\dfrac{55+7\sqrt{61}}{2}$}
\end{qoptions}
\qn{6} Let f(x) = $\tan^{-1}\!\left(\dfrac{2}{\text{x}}\right)$,g(x) = $\tan^{-1}\dfrac{2}{|\text{x}|}$ \quad x $\neq$ 0 and h(x) = tan($\cos^{-1}$(sinx)) then
\begin{qoptions}
\opt{A}{h($f$(x)) + h(g(x)) = $\dfrac{4}{\text{x}}$ $\forall$ x $>$ 0}
\opt{B}{h($f$(x)) + h(g(x)) = x \ $\forall$ x $>$ 0}
\opt{C}{h($f$(x)) + h(g(x)) = 0 \ $\forall$ x $<$ 0}
\opt{D}{2h($f$(x)) + 4h(g(x)) = 3x \ $\forall$ x $>$ 0}
\end{qoptions}
\sect{SECTION-I (ii)}
\paraq{7}{8}
Consider a plane P whose equation is x -- y = 0 and two points A(3,2,4) and B(7,0,--10)\par
\qn{7} Coordinates of point P, which lies on plane $\Pi$ such that |PA + PB| is minimum, are-
\begin{qoptions}
\opt{A}{(21,21,9)}
\opt{B}{$\left(\dfrac{21}{4},\dfrac{21}{4},\dfrac{9}{4}\right)$}
\opt{C}{$\left(\dfrac{21}{8},\dfrac{21}{8},\dfrac{9}{4}\right)$}
\opt{D}{$\left(\dfrac{21}{16},\dfrac{21}{16},\dfrac{9}{16}\right)$}
\end{qoptions}
\qn{8} Angle between $\overrightarrow{\text{AB}}$ and vector perpendicular to plane $\Pi$ is-
\begin{qoptions}
\opt{A}{$\cos^{-1}\!\left(\dfrac{1}{2\sqrt{3}}\right)$}
\opt{B}{$\dfrac{\pi}{4}$}
\opt{C}{$\cos^{-1}\!\left(\dfrac{1}{6}\right)$}
\opt{D}{$\cos^{-1}\!\left(\dfrac{1}{12}\right)$}
\end{qoptions}
\paraq{9}{10}
Let A,B,C be 3 sets such that\par
A = \{(x, y) : x = 2 cos $\theta$, y = 2 sin $\theta$, $\theta\in$ R\}\par
B = $\left\{(\text{x, y}):\dfrac{\text{x}+1}{\cos 2205^\circ}=\dfrac{\text{y}}{\sin 2205^\circ}=\text{r, r}\in\text{R}\right\}$\par
C = $\left\{(\text{x, y}):\left(\text{x}-\sqrt{3}\right)^2+(\text{y}-1)^2\leqslant\text{a}^2;\ \text{a}\in\text{R}\right\}$\par
Let C$_1$,C$_2$ be the curves represented by A and B respectively while R be the region represented by C.\par
\qn{9} A points P(x$_1$,y$_1$) and Q(x$_2$,y$_2$) be such that P,Q $\in$ A $\cap$ B, then point of intersection of tangents drawn to C$_1$ at P and Q is
\begin{qoptions}
\opt{A}{(2,3)}
\opt{B}{(--4,4)}
\opt{C}{(3,--3)}
\opt{D}{(--5,5)}
\end{qoptions}
\qn{10} If minimum value of `a' for which A $\cap$ C = A is $\alpha$, then length of curve C$_2$ lying in R corresponding to a = $\alpha$, is
\begin{qoptions}
\opt{A}{$\sqrt{50}$}
\opt{B}{$\sqrt{52}$}
\opt{C}{$\sqrt{58}$}
\opt{D}{$\sqrt{57}$}
\end{qoptions}
\sect{SECTION-II}
\qn{1} Let $\displaystyle\lim_{\text{x}\to\infty}$ x ln$\left(\text{e}\left(1+\dfrac{1}{\text{x}}\right)^{1-\text{x}}\right)$ equals $\dfrac{\text{m}}{\text{n}}$ where m and n are relatively prime positive integer. Find 2m -- n.
\qn{2} Let t$_{100}$ = $\displaystyle\sum_{\text{r}=1}^{100}\dfrac{1}{\left({}^{100}\text{C}_\text{r}\right)^5}$ and S$_{100}$ = $\displaystyle\sum_{\text{r}=0}^{100}\dfrac{\text{r}}{\left({}^{100}\text{C}_\text{r}\right)^5}$, then the value of $\sec\!\left(\cos^{-1}\!\left(\dfrac{\text{S}_{100}-50}{150\text{t}_{100}}\right)\right)$ equals
\qn{3} A curve passes through (1, 5/6) and the slope of tangent at any point (x, y) is x$^2$ -- 3x $\forall$ x $\in$ R. The point of minimum ordinate on the curve where x $>$ 0 is (a, b), then find the value of 2(a + b).
\qn{4} If reciprocal of radius of the smaller circle which touches all the lines represented by the equation x$^2$y -- 2xy$^2$ -- 4xy = 0 is k$\left(\cos\dfrac{\pi}{\text{p}}\right)^2$, then find the value of (p -- k).
\qn{5} Let f be monic cubic polynomial such that f(1) = 1$^4$ -- 1, f(2) = 2$^4$ -- 2 and f(3) = 3$^4$ -- 3. If f(4) = N, then find the number of prime factors of N.
\qn{6} The quadratic polynomial P(x) satisfy the equation P(x) -- P(x + 1) = 2x -- 1 $\forall$ x $\in$ R. If the greatest value of P(x) in ($-\infty$, 0] is 2, then value of P(--1) + P(1) is equal to
\qn{7} If the locus of middle point of point of contact of tangent drawn to the parabola y$^2$ = 8x and the foot of perpendicular drawn from its focus to the tangent is a conic then find the length of latus rectum of the conic.
\qn{8} Let $f$ be a differentiable function such that 1 $<$ $f'$(x) $<$ 2 $\forall$ x $\in$ R and $f$(3) = 1. If range of all possible values of $f$(5) lies in ($\lambda$,$\mu$), then value of $\dfrac{2\lambda}{\mu}$ is equal to
\raggedcolumns
\papersec{2}
\partsec{1}{PHYSICS}
\sect{SECTION-I (i)}
\qn{1} 1 mol He gas undergoes a cyclic process ABCA as shown in the figure. Process AB, BC \& CA are isothermal, isochoric and adiabatic respectively. Temperature of gas at point A is T$_0$. On the basis of given information identify correct statements. (2$^{2/3}$ = 1.5) and (ln2 = 0.7)\par
\qfig{p2_phy_s1_q01_1}
\begin{qoptions}
\opt{A}{Change in internal energy in process BC is --RT$_0$/2.}
\opt{B}{Work done by gas in process CA is --RT$_0$/2.}
\opt{C}{Work done by gas in cyclic process is RT$_0$/5.}
\opt{D}{Efficiency of cyclic process is 2/7.}
\end{qoptions}
\qn{2} A particle of mass m is thrown horizontally from highest point A of tower AB with speed 30m/sec. At the same time another particle of mass 2m is thrown vertically upward from point C on ground with same speed 30 m/sec. Both particles strike in air and stuck to each other after two seconds of projection. Combined mass fall on ground at point D. On the basis of this information which of the option given below are correct ? ($\sqrt{2}$ = 1.4)\par
\qfig{p2_phy_s1_q02_1}
\begin{qoptions}
\opt{A}{Height of tower AB and distance BC, both are 60 meter.}
\opt{B}{After collision combined mass move with speed 20m/sec.}
\opt{C}{Distance BD is equal to 88 meter.}
\opt{D}{Time of flight of particle A is 2$\sqrt{3}$ sec. (If there is no collision)}
\end{qoptions}
\qn{3} A particle of specific charge $\alpha$ is projected from origin at t = 0 with velocity $\vec{\text{v}}$ = v$_0\hat{\text{i}}$ + v$_0\hat{\text{k}}$ in magnetic field $\vec{\text{B}}$ = --B$_0\hat{\text{k}}$. Mass of the particle is 1 SI unit. On the basis of given information find correct options.
\begin{qoptions}
\opt{A}{Velocity of particle at t = $\dfrac{\pi}{\alpha\text{B}_0}$ is --v$_0$($\hat{\text{i}}$ -- $\hat{\text{k}}$).}
\opt{B}{Speed of the particle at t = $\dfrac{\pi}{\alpha\text{B}_0}$ is $\dfrac{\sqrt{5}}{2}$v$_0$.}
\opt{C}{Distance travelled by particle upto time t = $\dfrac{2\pi}{\alpha\text{B}_0}$ is $\dfrac{2\sqrt{2}\text{v}_0\pi}{\alpha\text{B}_0}$.}
\opt{D}{Displacement of particle upto time t = $\dfrac{3\pi}{\alpha\text{B}_0}$ is $\dfrac{3\pi\text{v}_0}{\alpha\text{B}_0}$.}
\end{qoptions}
\qn{4} During fair weather conditions atmosphere of the Earth can be considered as a good conductor with resistivity $\rho$ depending on the height h from the surface of the Earth as $\dfrac{1}{\rho}$ = a + bh$^2$, where a and b are constants. It is also known that near the surface of the Earth is a downward electric field E$_0$, which causes a radial current in the atmosphere. Assume E$_0$ is constant in magnitude.\par
\qfig{p2_phy_s1_q04_1}
\begin{qoptions}
\opt{A}{The current density J decreases with increase in height.}
\opt{B}{The electrons will be travelling in radially outward direction.}
\opt{C}{The charge density will be negative at any general height.}
\opt{D}{The charge density at height h is $\dfrac{2\varepsilon_0\text{E}_0}{\text{R}+\text{h}}$ in magnitude.}
\end{qoptions}
\qn{5} On a smooth horizontal table there are two blocks and two thin light discs. The axis of one of them is fixed, so that that disc cannot move, but can rotate about its axis. A piece of light inextensible thread covers the pulleys, to the ends of the thread we attach identical loads of mass M each. The threads are always taut. A constant force F is exerted on the movable disc as shown in the figure (top view). The thread does not slip relative to the discs.\par
\qfig{p2_phy_s1_q05_1}
\begin{qoptions}
\opt{A}{The left disc does not rotate but the right disc rotates in the clockwise direction.}
\opt{B}{The left disc rotates in clockwise direction and the right disc rotates in anticlockwise direction.}
\opt{C}{The acceleration of both the blocks is F/2M in opposite directions.}
\opt{D}{The net torque on the fixed disc is zero (or negligibly small).}
\end{qoptions}
\qn{6} Unpolarized light is incident on a prism of angle 45$^\circ$ and refractive index 4/3. The ray partially reflected at the first surface is completely polarized. Choose the correct statement(s).
\begin{qoptions}
\opt{A}{The angle of deviation is approximately 18.7$^\circ$.}
\opt{B}{The transmitted ray at the second face is completely polarized.}
\opt{C}{The transmitted ray at the first face is completely polarized.}
\opt{D}{There is no partial reflection at the second face.}
\end{qoptions}
\sect{SECTION-I (ii)}
\qn{7} A round object is under pure rolling on a rough surface. Different condition of motion of round object on different rough surfaces are shown by List-I. In all cases, there is no slipping at point of contact. Direction of velocity of centre of object is forward direction and opposite to velocity of centre is reverse direction. Now match List-I with List-II.\par
{\scriptsize\setlength{\tabcolsep}{1.5pt}\noindent\begin{tabular}{|L{.65cm}|L{4.4cm}|L{.55cm}|L{1.6cm}|}
\hline
\multicolumn{2}{|c|}{List-I} & \multicolumn{2}{c|}{List-II} \\ \hline
(I) & \qfigi{p2_phy_s2_q07_1} & (P) & Friction is in forward direction. \\ \hline
(II) & \qfigi{p2_phy_s2_q07_2} & (Q) & Friction is in reverse direction. \\ \hline
(III) & \qfigi{p2_phy_s2_q07_3} & (R) & Angular momentum of object is conserved with respect to point of contact of object. \\ \hline
(IV) & \qfigi{p2_phy_s2_q07_4} & (S) & Friction is zero. \\ \hline
 & & (T) & Angular momentum is conserved with respect to centre of object \\ \hline
\end{tabular}}\par
\begin{qoptions}
\opt{A}{I $\rightarrow$ P;II $\rightarrow$ Q;III $\rightarrow$ S,T;IV $\rightarrow$ Q,R}
\opt{B}{I $\rightarrow$ P,Q;II $\rightarrow$ Q,T;III $\rightarrow$ S,T;IV $\rightarrow$ Q,R}
\opt{C}{I $\rightarrow$ P;II $\rightarrow$ Q,S;III $\rightarrow$ R,T;IV $\rightarrow$ Q,R}
\opt{D}{I $\rightarrow$ P,Q;II $\rightarrow$ Q,T;III $\rightarrow$ R,T;IV $\rightarrow$ Q,T}
\end{qoptions}
\qn{8} A capillary is partially dipped in a liquid. Angle of contact between liquid and material of capillary is zero. Length of capillary above liquid is $\ell_1$ while length of capillary inside liquid is $\ell_2$. Liquid rises upto a light h in capillary. Radius of capillary is negligible with respect to $\ell_1$, $\ell_2$ \& h.\par
\qfig{p2_phy_s2_q08_1}
{\footnotesize\setlength{\tabcolsep}{2pt}\noindent\begin{tabular}{|L{.6cm}|L{3.3cm}|L{.6cm}|L{2.2cm}|}
\hline
\multicolumn{2}{|c|}{List-I} & \multicolumn{2}{c|}{List-II} \\ \hline
(I) & If $\ell_1$ $>$ h \& $\ell_2$ $>$ h and capillary is taken out of liquid completely. & (P) & Both upper and lower free surface of liquid in capillary are hemispherical. \\ \hline
(II) & If $\ell_1$ $>$ h, $\ell_2$ $<$ h and capillary is taken out of liquid completely. & (Q) & Length of liquid in capillary is 2h. \\ \hline
(III) & $\ell_1$ $>$ h and $\ell_2$ $>$ h. If system (beaker + capillary) is placed inside International Space Station (ISS) where gravity is not present but atmospheric pressure is present. Now capillary is taken out of liquid ($\ell_1$, $\ell_2$ and h are measured when experiment is conducted on earth). & (R) & Length of liquid in capillary is $\ell_1$ + $\ell_2$. \\ \hline
(IV) & If $\ell_1$ $<$ h, $\ell_2$ $<$ h and capillary is taken out of liquid & (S) & Both upper and lower free surface of liquid in capillary are flat. \\ \hline
 & & (T) & Lower free surface of liquid in capillary is not hemispherical. \\ \hline
\end{tabular}}\par
\begin{qoptions}
\opt{A}{I $\rightarrow$ P,Q;II $\rightarrow$ R,T;III $\rightarrow$ R;IV $\rightarrow$ R,T}
\opt{B}{I $\rightarrow$ P,R;II $\rightarrow$ P,Q;III $\rightarrow$ R,S;IV $\rightarrow$ R,T}
\opt{C}{I $\rightarrow$ P,Q;II $\rightarrow$ T;III $\rightarrow$ R;IV $\rightarrow$ R,T}
\opt{D}{I $\rightarrow$ P,R;II $\rightarrow$ P,R;III $\rightarrow$ Q,R;IV $\rightarrow$ Q,T}
\end{qoptions}
\qn{9} Three metallic bars 1, 2, 3 are arranged as shown is figure, with total number of free charge carriers in ratio N$_1$ : N$_2$ : N$_3$ = 1 : 3 : 2; resistivity ratio $\rho_1$: $\rho_2$ : $\rho_3$ = 2 : 1 : 3; lengths in ratio $\ell_1$ : $\ell_2$ : $\ell_3$ $\equiv$ 2 : 2 : 3 for 1, 2, and 3 bars respectively (radius of cross-section shown in figure), carry current i as shown, then match the columns.\par
\qfig{p2_phy_s2_q09_1}
{\footnotesize\setlength{\tabcolsep}{2pt}\noindent\begin{tabular}{|L{.6cm}|L{3.8cm}|L{.6cm}|L{1.3cm}|}
\hline
\multicolumn{2}{|c|}{List-I} & \multicolumn{2}{c|}{List-II} \\ \hline
(I) & If P$_1$, P$_2$ and P$_3$ are power dissipated across AB, BC, and CD respectively then $\dfrac{\text{P}_2\text{P}_3}{\text{P}_1^2}$ & (P) & $\dfrac{81}{32}$ \\ \hline
(II) & If E$_1$, E$_2$ and E$_3$ are magnitude of Electric-fields across AB, BC and CD respectively then $\dfrac{\text{E}_3^2}{\text{E}_1.\text{E}_2}$ & (Q) & $\dfrac{1}{32}$ \\ \hline
(III) & If $\upsilon_{\text{d}_1}$, $\upsilon_{\text{d}_2}$ and $\upsilon_{\text{d}_3}$ are drift speeds in bar AB, BC and CD respectively then $\dfrac{\upsilon_{\text{d}_2}\upsilon_{\text{d}_3}}{\upsilon_{\text{d}_1}^{\,2}}$ & (R) & $\dfrac{1}{4}$ \\ \hline
(IV) & If V$_1$,V$_2$, V$_3$ are potential differences across AB, BC and CD respectively then $\dfrac{\text{V}_2.\text{V}_3}{\text{V}_1^{\,2}}$ & (S) & $<$ 1 \\ \hline
 & & (T) & $\dfrac{5}{32}$ \\ \hline
\end{tabular}}\par
\begin{qoptions}
\opt{A}{I $\rightarrow$ Q,S;II $\rightarrow$ P;III $\rightarrow$ Q,S;IV $\rightarrow$ R,S}
\opt{B}{I $\rightarrow$ Q,S;II $\rightarrow$ P;III $\rightarrow$ R,S;IV $\rightarrow$ Q,S}
\opt{C}{I $\rightarrow$ P,S;II $\rightarrow$ P,S;III $\rightarrow$ Q,S;IV $\rightarrow$ R,S}
\opt{D}{I $\rightarrow$ P,S;II $\rightarrow$ Q,S;III $\rightarrow$ R,S;IV $\rightarrow$ Q,T}
\end{qoptions}
\qn{10} In the circuit shown S$_1$ and S$_2$ are switches which remains closed for a long time. Now S$_3$ and S$_4$ are closed simultaneously. Just after S$_3$ and S$_4$ are closed, I$_1$ and I$_2$ are the currents as shown in circuit. In the List-l some parameters are asked just after S$_3$ and S$_4$ are closed. Match these parameters with the answers given in List-II ($\varepsilon$ = 10v, R = $\dfrac{8}{3}\Omega$).\par
\qfig{p2_phy_s2_q10_1}
{\footnotesize\setlength{\tabcolsep}{2pt}\noindent\begin{tabular}{|L{.6cm}|L{3.5cm}|L{.6cm}|L{1.7cm}|}
\hline
\multicolumn{2}{|c|}{List-I} & \multicolumn{2}{c|}{List-II} \\ \hline
(I) & I$_1$ (in amperes) is & (P) & 2.5 \\ \hline
(II) & I$_2$ (in amperes) is & (Q) & 50 \\ \hline
(III) & Rate of change in energy of capacitor is (in J/s) & (R) & Zero \\ \hline
(IV) & Rate of change of magnetic field inside the inductor (in T/s) (assume inductor is a solenoid having 1000 turns per unit length and its volume equal to $\dfrac{16}{3}\times10^{-3}$ m$^3$) & (S) & Positive \\ \hline
 & & (T) & Negative \\ \hline
\end{tabular}}\par
\begin{qoptions}
\opt{A}{I $\rightarrow$ P,S;II $\rightarrow$ P,T;III $\rightarrow$ Q,T;IV $\rightarrow$ P,S}
\opt{B}{I $\rightarrow$ P,S;II $\rightarrow$ Q,T;III $\rightarrow$ Q,T;IV $\rightarrow$ P,S}
\opt{C}{I $\rightarrow$ P,Q;II $\rightarrow$ P,T;III $\rightarrow$ P,S;IV $\rightarrow$ R,S}
\opt{D}{I $\rightarrow$ P,R;II $\rightarrow$ P,T;III $\rightarrow$ P,S;IV $\rightarrow$ P,T}
\end{qoptions}
\sect{SECTION-II}
\qn{1} A resistance R and nonlinear device D are connected with an ideal battery as shown in the given network. Current I through the non linear device varies with the voltage V according to the equation I$^2$ = V$\cdot$K where K is positive constant. If current in the circuit is given by equation I = $\dfrac{-\text{RK}+\sqrt{\text{R}^2\text{K}^2+\text{xV}_0\text{K}}}{2}$ then find x.\par
\qfig{p2_phy_s3_q01_1}
\qn{2} An equiconvex lens of focal length 20cm is placed in air medium. It is divided into two parts about an axis $\perp$ to principal axis. One part is inverted about axis of division of lens. Both parts are in contact with each other. Space between lenses is filled by a liquid of unknown refractive index. When light rays parallel to principal axis fall on this combination no deflection is produced in light rays by this combination. Find refractive index of liquid. \{Refractive index of lens is $\dfrac{3}{2}$\}\par
\qfig{p2_phy_s3_q02_1}
\qn{3} A light beam of wavelength 400 nm is incident on a metal of work function 2.2 eV. A particular e$^-$ absorb this photon and make two collision before coming out of metal. Assuming that 10\% of existing energy is lost to the metal in each collision find final KE of this electron in eV as it comes out of metal. (hc = 12400 eV\AA)
\qn{4} Given figure shows a battery connected with various resistances and inductances through a switch. In the begining switch is open. Values of various parameters are\par
\qfig{p2_phy_s3_q04_1}
E = 15 volt, R$_1$ = 30$\Omega$\par
R$_2$ = 10$\Omega$, R$_3$ = 20$\Omega$\par
and L = 3 Henry\par
At t = 0 switch is closed. The graph which shows variation of current flowing through battery is given by\par
\qfig{p2_phy_s3_q04_2}
Find I$_2$ -- I$_1$ (in ampere)
\qn{5} A particle is moving on x-axis and has potential energy U = 5 -- 20x + 2x$^2$ joule where x is position of particle in meter. Particle is released at x = --2 meter. If mass of particle is 2 kg then maximum velocity of the particle is k$\sqrt{2}$ m/s. Find k.
\qn{6} Two planets of mass M \& 4M are shown in the figure. Radius of the planets R and 2R respectively and distance between centres of the planets in 6R. A small particle is thrown from planet of mass M along the line joining centres of both particles. Minimum speed of the particle so that it can reach other planet is given by $\sqrt{\dfrac{\text{x}\cdot\text{GM}}{5\text{R}}}$. Find x.\par
\qfig{p2_phy_s3_q06_1}
\qn{7} Four thin metallic plates have rectangular shapes with dimensions L $\times$ a or L $\times$ 2a. All those plates are placed parallel to each other with equal spacing d between consecutive neighboring plate as shown at the picture below. Top and bottom plates are connected with wire. Determine effective capacitance C$_{\text{AB}}$ measured between contacts A and B connected to the middle plates. If C$_{\text{AB}}$ = $\dfrac{\text{x}\varepsilon_0\text{aL}}{\text{d}}$, find x.\par
\qfig{p2_phy_s3_q07_1}
\qn{8} YDSE is carried out in a liquid of refractive index $\mu$ = 3/2. A thin film of air is formed in front of the lower slit as shown in the figure. A maxima is formed at point O which is mid point of screen. If thickness of air film is 4 $\mu$m then find order of maxima at point O. Wavelength of light in air is 4000\AA.\par
\qfig{p2_phy_s3_q08_1}
\raggedcolumns
\papersec{2}
\partsec{2}{CHEMISTRY}
\sect{SECTION-I (i)}
\qn{1} Which of the following set(s) of complexes is/are correct with respect to the properties indicated against each?
\begin{qoptions}
\opt{A}{[V(CO)$_6$], K$_4$[Mn(CN)$_6$], Cs$_2$[Ni$^{4+}$F$_6$] - Low spin octahedral complex with one unpaired electron.}
\opt{B}{K$_2$[NiCl$_4$], [NiCl(PPh$_3$)$_3$], [NiCl$_2$(PPh$_3$)$_2$] - Tetrahedral geometry containing two unpaired electrons.}
\opt{C}{Cs[AgF$_4$], [Cu(NH$_3$)$_4$] (NO$_3$)$_2$, K$_2$[Zn(CN)$_4$] - Square planar geometry and diamagnetic.}
\opt{D}{Na[Co(CO)$_4$], [Ni(CO)$_4$], Na$_2$[Fe(CO)$_4$] - Tetrahedral geometry and diamagnetic.}
\end{qoptions}
\qn{2} Which of the following anion is uninegative bidentate ligand?
\begin{qoptions}
\opt{A}{Acetylacetonate ion}
\opt{B}{Oxalate ion}
\opt{C}{Dimethylglyoximate ion}
\opt{D}{None of these}
\end{qoptions}
\qn{3} In which of the following reaction gives mixture of benzne and dibasic acid as product?
\begin{qoptions}
\opt{A}{\qfigi{p2_chem_s1_q03_A}}
\opt{B}{\qfigi{p2_chem_s1_q03_B}}
\opt{C}{\qfigi{p2_chem_s1_q03_C}}
\opt{D}{\qfigi{p2_chem_s1_q03_D}}
\end{qoptions}
\qn{4} Which of the following reaction(s) is/are related to extraction of wrought iron from haematite ore:
\begin{qoptions}
\opt{A}{MnO + SiO$_2$ $\xrightarrow{\Delta}$ MnSiO$_3$}
\opt{B}{3Fe$_2$O$_3$(s) + CO(g) $\xrightarrow{\Delta}$ 2Fe$_3$O$_4$(s) + CO$_2$(g)}
\opt{C}{CaO + P$_4$O$_{10}$ $\xrightarrow{\Delta}$ Ca$_3$(PO$_4$)$_2$}
\opt{D}{FeCO$_3$ $\xrightarrow{\Delta}$ FeO + CO$_2\uparrow$}
\end{qoptions}
\qn{5} For 2 moles of a diatomic gas undergoing a polytropic process for which :\par
$\dfrac{\text{dP}}{\text{dV}} = \dfrac{\text{P}^2}{2\text{RT}}$. (symbols have usual meaning)\par
Choose the correct option(s).
\begin{qoptions}
\opt{A}{Molar heat capacity for the process = 3R}
\opt{B}{Plot between `P' and `V' for the process is linear}
\opt{C}{Plot between `P' and `V$^{2}$' for the process is linear}
\opt{D}{Plot between `T' and `V$^{2}$' for the process is linear}
\end{qoptions}
\qn{6} \qfigi{p2_chem_s1_q06_1}\par
Major products (P) and (Q) are :
\begin{qoptions}
\opt{A}{Product (P) = \qfigi{p2_chem_s1_q06_2} ; Product (Q) = \qfigi{p2_chem_s1_q06_3}}
\opt{B}{Product (P) =\par\qfigi{p2_chem_s1_q06_4}}
\opt{C}{Product (P) = \qfigi{p2_chem_s1_q06_5}\par Product (Q) = \qfigi{p2_chem_s1_q06_6}}
\opt{D}{Product (P) =\par\qfigi{p2_chem_s1_q06_7}}
\end{qoptions}
\sect{SECTION-I (ii)}
\qn{7}\par
{\scriptsize\setlength{\tabcolsep}{1pt}\noindent\begin{tabular}{|L{.6cm}|L{5.2cm}|L{.6cm}|L{1.5cm}|}
\hline
 & \multicolumn{1}{c|}{List-I (Reactions)} & & \multicolumn{1}{L{1.5cm}|}{List-II (A and B is differentiated by)} \\ \hline
(P) & Ph -- CH = O $\xrightarrow[\text{(2)H}^{+}]{\text{(1)Conc.HO}^{-}}$ A + B & (1) & Tollen's test \\ \hline
(Q) & CH$_3$ -- C $\equiv$ CH $\xrightarrow[\text{dil.H}_2\text{SO}_4]{\text{Hg}^{+2}}$ A\par CH$_3$ -- C $\equiv$ CH $\xrightarrow[\text{(2) H}_2\text{O}_2\text{, HO}^{-}]{\text{(1) BH}_3\text{, THF}}$ B & (2) & Fehling test \\ \hline
(R) & \qfigi{p2_chem_s2_q07_4} & (3) & NaHCO$_3$ \\ \hline
(S) & \qfigi{p2_chem_s2_q07_5} & (4) & Iodoform test \\ \hline
 & & (5) & FeCl$_3$ \\ \hline
\end{tabular}}\par
\begin{qoptions}
\opt{A}{P $\rightarrow$ 3,4;Q $\rightarrow$ 1,2,4;R $\rightarrow$ 4;S $\rightarrow$ 5}
\opt{B}{P $\rightarrow$ 2;Q $\rightarrow$ 1,2;R $\rightarrow$ 3,4;S $\rightarrow$ 1,5}
\opt{C}{P $\rightarrow$ 2,5;Q $\rightarrow$ 2,4;R $\rightarrow$ 1,2;S $\rightarrow$ 3,4,}
\opt{D}{P $\rightarrow$ 1,5;Q $\rightarrow$ 2,3,4;R $\rightarrow$ 2,5;S $\rightarrow$ 1,2}
\end{qoptions}
\qn{8}\par
{\scriptsize\setlength{\tabcolsep}{1.5pt}\noindent\begin{tabular}{|L{.6cm}|L{3cm}|L{.6cm}|L{3.2cm}|}
\hline
 & List-I\par (Reagent) & & List-II\par (Observable change with given compounds) \\ \hline
(P) & Na$_2$O$_2$ solution & (1) & FeCl$_3$ solution \\ \hline
(Q) & Pb(CH$_3$COO)$_2$ solution & (2) & (NH$_4$)$_2$S solution \\ \hline
(R) & Na$_2$CO$_3$ solution & (3) & CuSO$_4$ solution \\ \hline
(S) & IO$_3^-$(aq) + H$^+$(aq) & (4) & KI solution \\ \hline
 & & (5) & Cr$_2$(SO$_4$)$_3$ solution \\ \hline
\end{tabular}}\par
\begin{qoptions}
\opt{A}{P $\rightarrow$ 2,3;Q $\rightarrow$ 1,2,3;R $\rightarrow$ 3,5;S $\rightarrow$ 2,3}
\opt{B}{P $\rightarrow$ 1,2,3,4,5;Q $\rightarrow$ 1,2,3,4,5;R $\rightarrow$ 1,2,3,5;S $\rightarrow$ 2,4}
\opt{C}{P $\rightarrow$ 3,4,5;Q $\rightarrow$ 2,3;R $\rightarrow$ 1,2;S $\rightarrow$ 1}
\opt{D}{P $\rightarrow$ 3,4;Q $\rightarrow$ 2,3;R $\rightarrow$ 1,2;S $\rightarrow$ 4,5}
\end{qoptions}
\qn{9}\par
{\scriptsize\setlength{\tabcolsep}{1.5pt}\noindent\begin{tabular}{|L{.6cm}|L{3.1cm}|L{.6cm}|L{3.4cm}|}
\hline
 & List - I\par Order & & List - II\par Correct Parameter \\ \hline
(P) & \qfigi{p2_chem_s2_q09_1} & (1) & Reactivity towards an electrophile or Lewis acid. \\ \hline
(Q) & \qfigi{p2_chem_s2_q09_2} & (2) & Resonance Energy \\ \hline
(R) & \qfigi{p2_chem_s2_q09_3} & (3) & At least one of the compound is aromatic \\ \hline
(S) & \qfigi{p2_chem_s2_q09_4} & (4) & Enthalpy of hydrogenation \\ \hline
\end{tabular}}\par
\begin{qoptions}
\opt{A}{P $\rightarrow$ 2,3;Q $\rightarrow$ 3;R $\rightarrow$ 1,2,3;S $\rightarrow$ 1,2,3}
\opt{B}{P $\rightarrow$ 3,4;Q $\rightarrow$ 1;R $\rightarrow$ 1,4;S $\rightarrow$ 2,3,4}
\opt{C}{P $\rightarrow$ 2;Q $\rightarrow$ 1,2;R $\rightarrow$ 3,4;S $\rightarrow$ 2}
\opt{D}{P $\rightarrow$ 1,4;Q $\rightarrow$ 3;R $\rightarrow$ 2,3;S $\rightarrow$ 2,3}
\end{qoptions}
\qn{10}\par
{\scriptsize\setlength{\tabcolsep}{1.5pt}\noindent\begin{tabular}{|L{.6cm}|L{2cm}|L{.6cm}|L{4cm}|}
\hline
 & List-I (species) & & List-II (Correct properties) \\ \hline
(P) & Na[BF$_4$] & (1) & Ionic bond \\ \hline
(Q) & CO & (2) & Covalent bond \\ \hline
(R) & O$_3$ & (3) & $\sigma$-dative bond \\ \hline
(S) & SiF$_4$ & (4) & $\pi$-dative bond \\ \hline
 & & (5) & All the atoms follow octet rule \\ \hline
\end{tabular}}\par
\begin{qoptions}
\opt{A}{P $\rightarrow$ 1,2,3,5;Q $\rightarrow$ 2,4,5;R $\rightarrow$ 2,3,5;S $\rightarrow$ 2,5}
\opt{B}{P $\rightarrow$ 1,2;Q $\rightarrow$ 4,5;R $\rightarrow$ 2,4;S $\rightarrow$ 1,3}
\opt{C}{P $\rightarrow$ 2,5;Q $\rightarrow$ 2,4;R $\rightarrow$ 3,5;S $\rightarrow$ 1,2}
\opt{D}{P $\rightarrow$ 2,3;Q $\rightarrow$ 4,5;R $\rightarrow$ 1,4;S $\rightarrow$ 1,4,5}
\end{qoptions}
\sect{SECTION-II}
\qn{1} Maximum number of atoms that can lie in a plane in the stable configuration of S$_2$F$_{10}$ molecule.
\qn{2} 2 litres of a hard water sample having hardness due to CaSO$_4$ is passed through anion exchange resin [R(OH)$_2$] and water coming out of resin required 5 ml 0.001 M H$_2$SO$_4$ for neutralisation. Calculate hardness of this sample in terms of CaCO$_3$. [Atomic mass : Ca = 40, S = 32]
\qn{3} In the following reaction scheme\par
\qfig{p2_chem_s3_q03_1}
Find the number of unpaired electron(s) in the complex [C]
\qn{4} A galvanic cell was constructed by dipping a zinc electrode (anode) in 0.1 M Zn(NO$_3$)$_2$ and hydrogen electrode (cathode) in which 0.2 M HIO$_3$(aq) solution is present.\par
If EMF of cell at 25$^\circ$C is x volt then find the value of $\dfrac{100\text{x}}{15.2}$\par
Given : K$_{\text{a(HIO}_3\text{)}}$ = 0.10 and partial pressure of H$_2$(g) = 0.1 bar; E$^0_{\text{Zn|Zn}^{2+}}$ = 0.76 V
\qn{5} If an octahedral homoleptic complex [ML$_\text{x}$]$^{\text{n}\pm}$ does not exhibit geometrical isomerism but it exhibits optical isomerism, then find out the value of expression (P + Q)\par
Where, P = Denticity of a ligand\par
\hspace*{1.2cm}Q = Number of Ligands(x)
\qn{6} \qfigi{p2_chem_s3_q06_A}\hfill\qfigi{p2_chem_s3_q06_B}\par
\qfigi{p2_chem_s3_q06_C}\par
\qfigi{p2_chem_s3_q06_D}\par
\qfigi{p2_chem_s3_q06_5}\par
\qfigi{p2_chem_s3_q06_6}\par
Calculate the number of reducing sugar = x\par
Calculate the number of sugar in which molecular weight increases by 210, whenever it undergo acetylation reaction with acetic anhydride = y\par
So x -- y =
\qn{7} Find the number of molecular species having at least one bond angle equal to 90$^\circ$.\par
SF$_4$, PF$_2$CI$_3$, XeO$_6^{4-}$, IF$_7$, C$\ell$F5, C$\ell$F$_3$, XeF$_4$, IO$_2$F$_2^-$, XeO$_3$F$_2$
\qn{8} A student is interested in coating a metal surface of 675 cm$^2$ area with 0.004 mm thickness with silver by passing a current of 96.5 amp with 50\% current efficiency through a silver nitrate aqueous solution. If the time required to complete the coating is `x' seconds then calculate the value of $\dfrac{\text{x}}{10}$.\par
(Given: Density of Ag is 10 g/cm$^3$, Atomic mass of Ag = 108)
\raggedcolumns
\papersec{2}
\partsec{3}{MATHEMATICS}
\sect{SECTION-I (i)}
\qn{1} Let (a -- 1) (x$^2$ + $\sqrt{3}$x + 1)$^2$ -- (a + 1)(x$^4$ -- x$^2$ + 1) $\leq$ 0 $\forall$ x $\in$ R, then which of the following is/are correct?
\begin{qoptions}
\opt{A}{a$\in\left[-\dfrac{1}{\sqrt{3}},\dfrac{4}{\sqrt{3}}\right]$}
\opt{B}{Largest possible value of a is $\dfrac{2}{\sqrt{3}}$}
\opt{C}{Number of possible integral values of a is 3}
\opt{D}{Sum of all possible integral values of a is `0'}
\end{qoptions}
\qn{2} Let $f$(x) = $\begin{cases}\displaystyle\int_0^{\text{x}}(5+|1-\text{t}|)\text{dt}, & \text{x}>2\\ 5\text{x}+1, & \text{x}\leqslant2\end{cases}$, then
\begin{qoptions}
\opt{A}{$f$(x) is differentiable at x = 2}
\opt{B}{$f$(x) is continuous but not differentiable at x = 2}
\opt{C}{$f$(x) is monotonically increasing $\forall$ x $\in$ R}
\opt{D}{$f$(4) = 25}
\end{qoptions}
\qn{3} If the system of equation 2x + y + z = 0, x -- 2y + z = 0 and ax -- y + 2z = 0 has infinitely many solutions and $f$(x) be a continuous function such that $f$(3 + x) + $f$(x) = 4 $\forall$ x $\in$ R, then $\displaystyle\int_0^{-2\text{a}}f(\text{x})\text{dx}$ is equal to\par
[Note : [k] denotes greatest integer function less than or equal to k.]
\begin{qoptions}
\opt{A}{--12}
\opt{B}{--2a}
\opt{C}{$\displaystyle\int_{-5}^{\text{a}}[\text{x}]\text{dx}$}
\opt{D}{2a}
\end{qoptions}
\qn{4} If $f$(x) = $\displaystyle\int\dfrac{\tan^3\text{x}}{2+\tan^2\text{x}}\text{dx} = \ln\left|\dfrac{1+g(\text{x})}{\cos\text{x}}\right|$ + C, where $f$(0) = ln2 and C is the constant of integration, then\par
[Note : [k] denotes greatest integer function less than or equal to k.]
\begin{qoptions}
\opt{A}{$\displaystyle\lim_{\text{x}\to0}\dfrac{g(\text{x})}{\sqrt{\text{x}^2-\text{x}^2\cos\text{x}}}$ = 2}
\opt{B}{$\displaystyle\int_0^{\pi/2}g(\text{x})\text{dx} = \tan^{-1}\!\left(\dfrac{1}{2}\right)+\tan^{-1}\!\left(\dfrac{1}{3}\right)$}
\opt{C}{$\displaystyle\lim_{\text{x}\to0^+}[\text{x}^2-g(\text{x})]$ = --1}
\opt{D}{$\displaystyle\int_0^{13\pi/3}\sqrt{1-g(\text{x})}\,\text{dx} = \dfrac{17}{2}$}
\end{qoptions}
\qn{5} Let (a$_1$, a$_2$, a$_3$ ... a$_\text{n}$) be a permutation of (1, 2, 3 ..... n) for which a$_1$ $>$ a$_2$ $>$ a$_3$ $>$ ..... a$_{\frac{\text{n}}{2}}$ and a$_{\frac{\text{n}}{2}}$ $<$ a$_{\frac{\text{n}}{2}+1}$ $<$ a$_{\frac{\text{n}}{2}+2}$ $<$ .... $<$ a$_\text{n}$ for n as an even positive integer. Also a$_1$ $>$ a$_2$ $>$ a$_3$ ....... $>$ a$_{\frac{\text{n}+1}{2}}$ and a$_{\frac{\text{n}+1}{2}}$ $<$ a$_{\frac{\text{n}+1}{2}+1}$ $<$ .... $<$ a$_\text{n}$ for n as an odd positive integers. Let the total number of permutation of n be p(n) if 100 $<$ p(n) $<$ 500, then value of n is/are
\begin{qoptions}
\opt{A}{10}
\opt{B}{11}
\opt{C}{12}
\opt{D}{13}
\end{qoptions}
\qn{6} Given ax$^2$ + bx + c $\geq$ 0 ; bx$^2$ + cx + a $\geq$ 0 and cx$^2$ + ax + b $\geq$ 0, where a $\neq$ b $\neq$ c and a, b, c $\in$ R. Then $\dfrac{\text{a}^2+\text{b}^2+\text{c}^2}{\text{ab}+\text{bc}+\text{ca}}$ can not take value(s)
\begin{qoptions}
\opt{A}{$\dfrac{2}{3}$}
\opt{B}{$\dfrac{1}{2}$}
\opt{C}{$\dfrac{3}{2}$}
\opt{D}{$\dfrac{7}{2}$}
\end{qoptions}
\sect{SECTION-I (ii)}
\qn{7} Match the following\par
{\scriptsize\setlength{\tabcolsep}{1.5pt}\noindent\begin{tabular}{|L{.6cm}|L{5.6cm}|L{.6cm}|L{.8cm}|}
\hline
 & \multicolumn{1}{c|}{List-I} & & List-II \\ \hline
(P) & Let $f$(x) be a non negative continuous $\forall$ x $\geq$ 0, $f$(0) = 0.\par If $f'$(x) $\leq$ $f$(x) sin x $\forall$ x $\geq$ 0, then the value of $f\left(\dfrac{3\pi}{2}\right)$ is & (1) & 4 \\ \hline
(Q) & If a(x$^2$ + 1) + 12x + 6 -- 3x$^2$ $<$ $\sin^{-1}\sin(100)+\cos^{-1}\cos(100)$ $\forall$ x $\in$ R, a $\in$ I then a can be equal to & (2) & 3 \\ \hline
(R) & If two events A and B are such that P(A$^\text{C}$) = 0.35, P(B) = 0.4 and P(A $\cap$ B) = 0.15, then 20P$\left(\dfrac{\text{B}}{\text{A}\cup\text{B}^\text{C}}\right)$ is & (3) & 0 \\ \hline
(S) & The number of elements in the range of the function\par $f$(x) = $\sqrt{\cos^{-1}(4\text{x}^3-3\text{x})-\dfrac{\pi}{2}}-\dfrac{(2\text{x}+3)!}{\sqrt{\text{x}+2}}$ is/are & (4) & 5 \\ \hline
 & & (5) & 8 \\ \hline
\end{tabular}}\par
\begin{qoptions}
\opt{A}{P $\rightarrow$ 3;Q $\rightarrow$ 5;R $\rightarrow$ 2;S $\rightarrow$ 1}
\opt{B}{P $\rightarrow$ 5;Q $\rightarrow$ 3;R $\rightarrow$ 2;S $\rightarrow$ 1}
\opt{C}{P $\rightarrow$ 5;Q $\rightarrow$ 3;R $\rightarrow$ 1;S $\rightarrow$ 2}
\opt{D}{P $\rightarrow$ 3;Q $\rightarrow$ 5;R $\rightarrow$ 1;S $\rightarrow$ 2}
\end{qoptions}
\qn{8} Consider a system of liner equations 3x + y -- z = 0, x -- $\dfrac{\text{py}}{4}$ + z = 2 and 2x -- y + 2z = q where p, q $\in$ I and p, q $\in$ [1, 10], then identify the correct statement(s).\par
{\scriptsize\setlength{\tabcolsep}{1.5pt}\noindent\begin{tabular}{|L{.6cm}|L{5.2cm}|L{.6cm}|L{.8cm}|}
\hline
 & \multicolumn{1}{c|}{List-I} & & List-II \\ \hline
(P) & Number of ordered pairs (p, q) for which system of equation has unique solution is & (1) & 1 \\ \hline
(Q) & Number of ordered pairs (p, q) for which system of equation has no solution is & (2) & 9 \\ \hline
(R) & Number of ordered pairs (p, q) for which system of equation has infinite solutions is & (3) & 10 \\ \hline
(S) & Number of ordered pairs (p, q) for which system of equation has atleast one solution is & (4) & 90 \\ \hline
 & & (5) & 91 \\ \hline
\end{tabular}}\par
\begin{qoptions}
\opt{A}{P $\rightarrow$ 5;Q $\rightarrow$ 3;R $\rightarrow$ 2;S $\rightarrow$ 1}
\opt{B}{P $\rightarrow$ 4;Q $\rightarrow$ 2;R $\rightarrow$ 1;S $\rightarrow$ 5}
\opt{C}{P $\rightarrow$ 4;Q $\rightarrow$ 2;R $\rightarrow$ 5;S $\rightarrow$ 1}
\opt{D}{P $\rightarrow$ 3;Q $\rightarrow$ 4;R $\rightarrow$ 5;S $\rightarrow$ 2}
\end{qoptions}
\qn{9} A man has 3 coins A, B and C, A is unbiased, the probability that a head will result when B is tossed is 2/3, the probability that a head will result when C is tossed is 1/3. One of the coin chosen at random and is tossed 3 times, giving a total of two heads and one tail.\par
{\scriptsize\setlength{\tabcolsep}{1.5pt}\noindent\begin{tabular}{|L{.6cm}|L{5.2cm}|L{.6cm}|L{1.2cm}|}
\hline
 & \multicolumn{1}{c|}{List-I} & & List-II \\ \hline
(P) & The probability that the chosen coin was A & (1) & 32/75 \\ \hline
(Q) & The probability that the chosen coin was B & (2) & 1/2 \\ \hline
(R) & The probability that the fourth toss will result in head & (3) & 209/1296 \\ \hline
(S) & The probability that the fourth toss will result in tail & (4) & 9/25 \\ \hline
 & & (5) & 215/1296 \\ \hline
\end{tabular}}\par
\begin{qoptions}
\opt{A}{P $\rightarrow$ 4;Q $\rightarrow$ 1;R $\rightarrow$ 2;S $\rightarrow$ 3}
\opt{B}{P $\rightarrow$ 5;Q $\rightarrow$ 4;R $\rightarrow$ 2;S $\rightarrow$ 3}
\opt{C}{P $\rightarrow$ 3;Q $\rightarrow$ 2;R $\rightarrow$ 5;S $\rightarrow$ 1}
\opt{D}{P $\rightarrow$ 4;Q $\rightarrow$ 1;R $\rightarrow$ 2;S $\rightarrow$ 2}
\end{qoptions}
\qn{10} Consider the matrix A = $\begin{bmatrix}3&-4\\1&-1\end{bmatrix}$ and B = $\begin{bmatrix}\text{a}&6\\0&1\end{bmatrix}$ and Let P be any orthogonal matrix and Q = PAP$^\text{T}$ and R = P$^\text{T}$Q$^\text{K}$P also S = PBP$^\text{T}$ and T = P$^\text{T}$S$^\text{K}$P\par
{\scriptsize\setlength{\tabcolsep}{1.5pt}\noindent\begin{tabular}{|L{.6cm}|L{4.2cm}|L{.6cm}|L{2.3cm}|}
\hline
 & \multicolumn{1}{c|}{List-I} & & List-II \\ \hline
(P) & If we very K from 1 to `n', then first row first column elements of R will form & (1) & G.P. with common ratio a \\ \hline
(Q) & If we very k from 1 to n, then 2$^{\text{nd}}$ row 2$^{\text{nd}}$ column elements of R will form & (2) & A.P with common difference 2 \\ \hline
(R) & If we very k from 1 to n, then first row first column elements of T will form & (3) & G.P. with common ratio `6' \\ \hline
(S) & If we very k from 3 to n, then first row 2$^{\text{nd}}$ column elements of T will represent the sum of & (4) & A.P with common difference --2 \\ \hline
 & & (5) & A.P with common difference --7 \\ \hline
\end{tabular}}\par
\begin{qoptions}
\opt{A}{P $\rightarrow$ 5;Q $\rightarrow$ 3;R $\rightarrow$ 4;S $\rightarrow$ 2}
\opt{B}{P $\rightarrow$ 4;Q $\rightarrow$ 2;R $\rightarrow$ 3;S $\rightarrow$ 5}
\opt{C}{P $\rightarrow$ 2;Q $\rightarrow$ 4;R $\rightarrow$ 1;S $\rightarrow$ 1}
\opt{D}{P $\rightarrow$ 3;Q $\rightarrow$ 1;R $\rightarrow$ 2;S $\rightarrow$ 5}
\end{qoptions}
\sect{SECTION-II}
\qn{1} The value of integral $\displaystyle\int_{1/\text{e}}^{\cos\text{x}}\dfrac{\text{t}}{1+\text{t}^2}\text{dt}+\int_{1/\text{e}}^{\sec\text{x}}\dfrac{\text{dt}}{\text{t}(1+\text{t}^2)}$ is equal to \underline{\hspace{1.5cm}}
\qn{2} Six points (x$_\text{i}$, y$_\text{i}$), i = 1, 2, ...., 6 are taken on the circle x$^2$ + y$^2$ = 4 such that $\displaystyle\sum_{\text{i}=1}^{6}$ x$_\text{i}$ = 12 and $\displaystyle\sum_{\text{i}=1}^{6}$ y$_\text{i}$ = 8. The line segment joining orthocentre of a triangle made by any three points and the centroid of the triangle made by other three points passes through a fixed points (h, k), then h + k is\underline{\hspace{1.5cm}}
\qn{3} Let N be the number of arrangements of the letters of the word ``LOVEISLIFE'' in which no two alike letters are together and containing the word LOVE. Then N is
\qn{4} If the equation of the curve on reflection of the ellipse $\dfrac{(\text{x}-4)^2}{16}+\dfrac{(\text{y}-3)^2}{9}$ = 1 about the line x -- y -- 2 = 0 is k$_1$x$^2$ + 9y$^2$ -- 160x -- k$_2$y + 292 = 0, then (k$_1$ + k$_2$)/13 is equal to \underline{\hspace{1.5cm}}
\qn{5} Let P is a point on the plane ax + by + cz = d. A point Q($\alpha$, $\beta$, $\gamma$) is taken on the line OP such that OP. OQ = d$^2$, then $\dfrac{2\text{d}(\text{a}\alpha+\text{b}\beta+\text{c}\gamma)}{\alpha^2+\beta^2+\gamma^2}$ is equal to (where O is origin)
\qn{6} Die A has 2 white and 4 red faces whereas die B has 4 white and 2 red faces. A coin is flipped once. If it shows a head, the game continues by throwing die A, if it shows tail, then die B is to be used. If the probability that die A is used is $\dfrac{32}{33}$ when it is given that red turns up every time in first n throws. Then find the value of n \underline{\hspace{1.2cm}}.
\qn{7} Let A = $\begin{bmatrix}\text{x}^3+1&\text{x}+5&3\text{x}+2\\ \text{y}+1&-6\text{x}^2+2&\text{z}-1\\ 2&3&9\text{x}+6\end{bmatrix}$ and |A| = 3\par
If $f$(x) = tr(B$^{-1}$) and B = adj(A), then global minimum value of $f$(x) in x $\in$[0, 6] is p and it is achieved at `n' number of value(s) of x then n + p is equal to
\qn{8} Let the circle S : x$^2$ + y$^2$ -- ax -- by + c = 0 intersects the pair of straight lines xy -- 4x -- 3y + 12 = 0 orthogonally and the circle lies in the first quadrant. If S touches the circle S$_1$ : (x --3)$^2$ + y$^2$ = r$^2$, r $\in$ N, then find the sum of all possible values of r.
\end{multicols}
\end{document}
